\section{Adaptive boundary classification of a supplied normal surface}
\label{sec:normal-boundary}

Report 47 supplies a useful refinement of normal component counting:
distinguish closed components from components meeting the ambient boundary,
and separate each class by orientability. We now integrate this refinement
with the maintained finite-manifold validator, edge-point interval model,
common-multiplicity reduction and independent certificate replay. The
public option \texttt{classify\_boundary=True} returns these additional
counts. It defaults to false, preserving the previous result structure and
version-one/version-two proofs.

This is a local normal-surface operation. The input is still a supplied
finite triangulation and an admissible integral normal vector. No
triangulation-to-diagram provenance or normal-vector search is added.
In particular, separating closed and boundary-touching component counts
does not give the Euler characteristic of each component or decide whether
a disconnected supplied surface contains an essential disc. The existing
\texttt{compressing\_disk} predicate still requires the entire supplied
surface to be connected, with Euler characteristic one, one boundary
circle and nonzero mod-two boundary class.

\subsection{Mark boundary edges in the maintained orbit universe}

Let $Q$ denote the normal surface after dividing out any common coordinate
factor. The maintained surface interval system has one point for each
intersection with an ambient edge, ordered in blocks by canonical edge
representative. Its orbits are precisely the connected components of $Q$.
Let $A$ be the union of blocks belonging to ambient boundary edges. A normal
surface component meets the ambient boundary exactly when its orbit meets
$A$: every nonempty normal boundary curve contains boundary-edge points,
and such a point lies on the boundary of the proper normal surface.

Empty blocks are omitted and adjacent marked blocks are merged. There are
$O(t)$ blocks for $t$ tetrahedra, with binary endpoints; no list of normal
discs or individual edge intersections is created. For the doubled vector
$2Q$, each edge block doubles in length, so a marked interval $[a,b)$ becomes
$[2a,2b)$. This detail depends on the maintained coordinate ordering.
Appending a second copy of $A$ shifted by the old total universe length
would instead describe a different ordering, and is not the construction
used here.

For nonempty $A$, add the existing interval-incidence cone relations:
translations identify all points within each marked interval, and a star
joins one representative from every marked interval. These relations merge
all original orbits touching $A$ into one, and leave every disjoint orbit
unchanged. If $C$ is the original orbit count, $C_A$ the count after coning,
and $T$ the number of touched components, then
\[
 C_A=C-T+1,\qquad T=C-C_A+1.
\]
This is a counting reduction, not an operation on the actual embedded
surface. It adds $O(t)$ interval pairings and retains binary endpoints.
The empty marked set has $T=0$ without a cone query.

\subsection{Orientation covers refine the boundary partition}

Write $C_2$ for the number of components of $2Q$, and $T_2$ for the number
of its components meeting the ambient boundary. Since the supported
ambient manifold is orientable, doubling normal coordinates gives the
normal orientation cover. An orientable component has two lifts; a
nonorientable component has one connected orientable lift. Boundary
presence is preserved by this covering. Therefore, with subscripts $b$
denoting components with boundary,
\[
 T=O_b+N_b,\qquad T_2=2O_b+N_b,
 \qquad O_b=T_2-T,\quad N_b=2T-T_2.
\]
The already verified total counts $O=C_2-C$ and $N=2C-C_2$ then give
closed counts $O_c=O-O_b$ and $N_c=N-N_b$. The implementation checks
nonnegativity and containment in the total orientation classes. If $b$
is the number of boundary circles of $Q$, a nonempty boundary also
requires $1\le T\le\min(C,b)$.

For an original vector $X=gQ$, the sheet-pairing argument of
Section~\ref{sec:normal-multiplicity} applies independently to components
with and without boundary. Thus
\[
 O_b(X)=gO_b+\lfloor g/2\rfloor N_b,
 \qquad N_b(X)=(g\bmod2)N_b.
\]
All reported counts concern $X$. The original source fingerprint,
Euler characteristic, normal-disc count and boundary-cohomology witness
continue to concern $X$ as well.

An even multiplier yields a further useful deduction. If $g$ is even,
$X$ consists of $g/2$ copies of the orientable normal surface $2Q$.
Consequently
\[
 O_b(X)=(g/2)T_2,\qquad N_b(X)=0.
\]
Computing $T$ is unnecessary in this case. In particular, for a primitive
surface with both orientation types the algorithm can omit the cone query
on $Q$, compute only $T_2$, and still determine all six requested original
counts. It does not claim to have recovered the unknown partition of $Q$.

\subsection{Choose queries from already established progress}

The three base queries still compute $C$, $C_2$ and $b$. The adaptive
policy then uses exact deductions before scheduling any further query:
\begin{itemize}
\item If $b=0$, both touched counts are zero. If $C=1$ and $b>0$,
      they are $T=1$ and $T_2=C_2$.
\item If the marked union is the entire edge-point universe, every
      component touches the boundary: $T=C$ and $T_2=C_2$.
\item If $b=1$, exactly one component touches the boundary, so $T=1$.
\item If $C_2=2C$, every component is orientable and $T_2=2T$.
      If $C_2=C$, every component is nonorientable and $T_2=T$.
      Either case needs at most one cone query.
\item If both orientation types occur and $g$ is even, compute only the
      doubled cone count when the preceding deductions do not suffice.
      Otherwise compute any still-needed counts on $Q$ and $2Q$.
\end{itemize}
Thus the feature uses zero, one or two additional orbit queries. This is
an adaptive reduction in the information that must be computed, rather
than a heuristic choice of an unverified result. A private research
ablation requests both cones whenever the boundary is nonempty; it is
not a new public option.

Each cone system still has $O(t)$ pairings and logarithmic dependence on
the edge-point universe size. At most five interval-orbit queries therefore
preserve the polynomial local bound in the encoded triangulation and
normal coordinates. No quasipolynomial bound for finding a useful normal
vector follows from this operation. The gains measured below concern
complete calls on supplied vectors, including validation, rather than
repeated queries on a prevalidated geometry object.

\subsection{Proof version three and independent replay}

Opting into classification with certificate recording produces
\texttt{normal-surface-topology-v3}. It binds the full original input,
records an exact coordinate divisor $g\ge1$, retains all three base
orbit proofs and the original boundary-cohomology witness, and includes
zero to two cone proofs. The claimed topology includes the six new fields:
the two boundary-presence totals and their four orientability subclasses.
The empty surface uses $g=1$ and six zero counts. As before, the checker
accepts any valid common divisor, not only the producer's gcd.

The checker validates the original source, divides coordinates exactly,
reconstructs both interval systems and the marked blocks, and replays
every supplied cone proof. It derives omitted counts directly from verified
base counts and the reconstructed marked support. It does not call the
producer's query planner, orbit search or gcd selection. Redundant valid
cone proofs are accepted and checked. Pure orientability permits inference
in either direction, so a sound alternate proof choice also verifies.
A proof containing only the doubled touched count can suffice for even
$g$; changing that divisor to an odd value without supplying the missing
information is rejected.

The producer's \texttt{max\_cycles} allowance is shared across the base
and all additional queries. Exhaustion during a cone query returns an
inconclusive result with query diagnostics, without partial topology
claims or a certificate. The verifier's \texttt{max\_operations} counts
all supplied orbit events before replay, including redundant cone proofs.
Cancellation callbacks propagate through both routes. Version-one and
version-two formats retain their earlier meanings; the older checker
is not claimed to understand version three.

\subsection{Independent geometry audit and measurements}

The pinned baseline is commit \texttt{ce180ce1e640}, loaded directly from
Git. The native audit checks 548 supplied surfaces against both Regina 7.4
component decomposition and report 47's delivered component implementation.
Every new count agrees; all default results and certificates equal the
actual maintained baseline exactly, with default proofs replayed in both
versions. Direct and adaptive classification agree and independently
verify on every input. The audit includes 404 three-query cases, 95
four-query cases and 49 five-query cases.

The corpus reuses 270 layered-solid-torus and boundary-cap inputs, adds
all 64 small combinations of an interior vertex sphere, a boundary vertex
disc and a M\"obius band in a fixed four-tetrahedron subdivision, and
includes vertex normal surfaces and compatible combinations in a
nine-tetrahedron solid-torus connected sum with $\mathbb{RP}^3$.
That final manifold is deliberately not a knot exterior: it supplies a
closed one-sided projective plane for validating the general local
classifier, without asserting nonexistent knot-diagram provenance.
The audit contains 74 inputs with closed orientable components, 15 with
closed nonorientable components, 473 with orientable components meeting
the boundary and 98 with nonorientable components meeting the boundary.
These classes overlap. Raw data: \path{data/normal-boundary-native-audit.json}.

Eight focused tests additionally exercise 5,001-bit multiplicities,
proper common divisors, the unknown primitive partition in the even case,
missing and forged cone proofs, mutated topology claims, shared exhaustion
and replay limits, caller cancellation and empty inputs. Replay is also
tested with the producer's planner, gcd routine and orbit search disabled.
The full maintained suite passes 966 tests with Regina in 213.279 seconds.

\begin{center}\small
\begin{tabular}{@{}lrrrrr@{}}
\toprule
Input & queries D/A & direct ms & adaptive ms & count ratio & proof ratio \\
\midrule
meridian8 & 5/3 & 6.151 & 4.970 & 1.222 & 1.238 \\
meridian32 & 5/3 & 75.565 & 69.700 & 1.088 & 1.049 \\
meridian96 & 5/3 & 1279.093 & 1274.040 & 0.999 & 0.971 \\
parallel 5000 & 5/3 & 6.251 & 5.122 & 1.210 & 1.242 \\
closed & 3/3 & 0.698 & 0.696 & 1.003 & 1.030 \\
one boundary & 5/4 & 1.811 & 1.530 & 1.185 & 1.149 \\
pure orientable & 5/4 & 5.637 & 4.236 & 1.338 & 1.279 \\
mixed odd & 5/5 & 5.559 & 5.521 & 0.997 & 1.110 \\
mixed even & 5/4 & 5.550 & 4.435 & 1.259 & 1.243 \\
even 5000 & 5/4 & 7.536 & 7.700 & 1.245 & 1.211 \\
odd 5000 & 5/5 & 5.771 & 5.749 & 0.989 & 0.986 \\
closed one-sided & 5/5 & 13.854 & 14.040 & 0.994 & 0.998 \\
\bottomrule
\end{tabular}
\end{center}


The paired benchmark has twelve supplied-vector inputs, eight arms, five
measured rounds and one excluded warmup: 1,920 measured complete calls and
384 warmup calls. Some small cases are batched. Direct and adaptive arms
answer the same six-field classification question, with common multiplicity
reduction enabled on both sides. Separate certified arms include proof
production and fresh independent replay. Identical direct and adaptive
controls expose timing variation. Two further arms compare the actual old
and new default operation, which answers a smaller question and is not
used as the classification speedup baseline. Source preparation, manifold
validation, coordinate validation and all queries are inside each timer;
policy binding, serialization and fixture loading are outside.
The 99.926-second run retains every input, shuffled arm order, sample,
query statistic, event count and source hash in
\path{../fast/results/normal_boundary_20261008.json}.

For the eight-tetrahedron meridian, the query count falls from five to
three and the paired count ratio is 1.222. Production plus replay has
ratio 1.238, with median times 10.121 and 8.263 ms. The compact proof falls
from 33,609 to 20,613 bytes and from 501 to 278 orbit events. The
32-tetrahedron input has a smaller count ratio, 1.088. At 96 tetrahedra,
the three mandatory queries dominate: the count ratio is 0.999 and the
certified ratio is 0.971, despite reducing orbit events from 26,591 to
24,038. This run establishes no elapsed improvement for that large case;
its identical-control ratios are 0.959 and 0.978.

On the four-tetrahedron mixture of a sphere, a boundary disc and a
M\"obius band, odd primitive multiplicity needs both cones. Doubling that
vector permits the new even shortcut: five queries become four, 55 cycles
become 44, and 351 events become 270. The proof decreases from 23,322 to
18,080 bytes. The paired count and certified ratios are 1.259 and 1.243.
The purely orientable sphere-plus-two-discs input also needs only one cone,
with ratios 1.338 and 1.279. A sphere plus a single M\"obius band already
has $b=1$, eliminating the original-surface cone and yielding ratios 1.185
and 1.149.

Large binary multiplicities retain the same query savings after gcd
reduction. The even $2^{5000}$ mixed input has paired count ratio 1.245
and certified ratio 1.211. Its separate count medians are 7.536 and
7.700 ms: the median of paired ratios need not equal a ratio of separate
medians, particularly with timing shifts between rounds. These raw data
are retained without selecting favorable rounds. The corresponding odd
$2^{5000}+1$ case still requires both cones and has ratios 0.989 and
0.986. The closed one-sided mixed input likewise requires both cones and
has ratios 0.994 and 0.998. The closed-only input requires no cone under
either policy.

Across all inputs the identical direct ratios range from 0.959 to 1.031
and the identical adaptive ratios from 0.978 to 1.008. The smaller default
operation has old/new ratios from 0.986 to 1.023. The primitive mixed-odd
case has identical query counts and proof size yet a certified ratio of
1.110; this is timing variation, not evidence of a saved query. Exact
query, cycle and proof-size reductions support the adaptive mechanism;
the elapsed measurements support modest gains on selected supplied
surfaces, not a universal speedup or a general recognition bound.
